\documentclass[11pt,a4paper]{article}

\def\ncourse {Probability Theory}
\def\ntype {problems}

\makeatletter

% packages
\usepackage{amssymb,amsfonts,amsmath,calc,tikz,pgfplots,geometry,mathtools}
\usepackage{color}   % May be necessary if you want to color links
\usepackage[svgnames]{xcolor}
\usepackage[hidelinks]{hyperref}
\usepackage{forest}
\usepackage{commath} % ew
\usepackage{esvect} % \vv makes vectors
\usepackage{centernot} % for the comparison
\usepackage{esdiff}
\usepackage{esint}
\usepackage{amsthm}
\usepackage{fancyhdr}
\usepackage{bm}
\usepackage{witharrows}
\usepackage{bookmark}
\usepackage{tikz-cd}
\usepackage{quiver}
\usepackage{bbm}
\usepackage{textcomp}
\usepackage{float}
\usepackage{gensymb}
\usepackage{cleveref}
\usepackage{environ}
\usepackage{comment}
\usepackage{cancel}
\usepackage{xparse}

% Inevitable cs
\usepackage{listings}
\usepackage[]{algorithm2e}

% tikz libraries
\usetikzlibrary{positioning}
\usetikzlibrary{matrix}
\usetikzlibrary{arrows}
\usetikzlibrary{bending}
\usetikzlibrary{arrows.meta}
\usetikzlibrary{decorations.markings}
\usetikzlibrary{decorations.pathreplacing}
\usetikzlibrary{decorations.markings}
\usetikzlibrary{patterns}
\usetikzlibrary{backgrounds}

% Page style setup
\pagestyle{fancy}
\fancyhfoffset{0pt}
\renewcommand{\headwidth}{\textwidth}
\geometry{margin=1in}
\pgfplotsset{compat=1.18}
\setlength{\headheight}{14.6pt}
\addtolength{\topmargin}{-1.6pt}
\hypersetup{
    colorlinks=false,
    linktoc=section,
    linkcolor=black,
}

%% maketitle setup
\ifx \nauthor\undefined
  \def\nauthor{Yeheli Fomberg}
\else
\fi

\ifx \nemail\undefined
  \def\nemail{\href{mailto:yp@campus.technion.ac.il}
  {\texttt{<yp@campus.technion.ac.il>}}}
\else
\fi

\ifx \ncoursehead \undefined
  \def\ncoursehead{\ncourse}
\fi

\lhead{\emph{\nouppercase{\leftmark}}}
\ifx \nextra \undefined
  \rhead{
    \ifnum\thepage=1
    \else
      \ncoursehead
    \fi}
\else
  \rhead{
    \ifnum\thepage=1
    \else
      \ncoursehead \ (\nextra)
    \fi}
\fi

\def\notes{notes}
\def\problems{problems}

\ifx \ntype \undefined
  \def\ntype{notes}
\else
\fi

\ifx \ntype \notes
  \let\@real@maketitle\maketitle
  \renewcommand{\maketitle}{\@real@maketitle\begin{center}
  \begin{minipage}[c]{0.9\textwidth}\centering\footnotesize
  These notes are not endorsed by the lecturers.
  I have revised them outside lectures to incorporate supplementary
  explanations, clarifications, and material for fun.
  While I have strived for accuracy, any errors or misinterpretations 
  are most likely mine.
  \end{minipage}\end{center}}
\else
  \ifx \ntype \problems
    \let\@real@maketitle\maketitle
    \renewcommand{\maketitle}{\@real@maketitle\begin{center}
    \begin{minipage}[c]{0.9\textwidth}\centering\footnotesize
    These problems are mostly based on the problem sets I was given
    when taking the course.
    
    While I have strived for accuracy and completeness,
    some of the problems would not have been here without the help of
    friends.
    \end{minipage}\end{center}}
  \else
  \fi
\fi


% theorem environments
\theoremstyle{definition}
\newtheorem{definition}{Definition}[section]
\newtheorem{remark}{Remark}[section]
\newtheorem{example}{Example}[section]
\newtheorem{exercise}{Exercise}[section]
\newtheorem{paradox}{Paradox}[section]
\newtheorem{entry}{Entry}[section]
\newtheorem*{solution}{Solution}
\newtheorem*{question}{Question}
\newtheorem*{answer}{Answer}
\newtheorem*{problem}{Problem}
\theoremstyle{plain}
\newtheorem{theorem}{Theorem}[section]
\newtheorem{proposition}[theorem]{Proposition}
\newtheorem{lemma}[theorem]{Lemma}
\newtheorem{corollary}[theorem]{Corollary}
\newenvironment{proofof}[1]{%
    \renewcommand{\proofname}{Proof of #1}%
    \begin{proof}
}{%
    \end{proof}
}
% tikz customization
\pgfarrowsdeclarecombine{twolatex'}{twolatex'}{latex'}{latex'}{latex'}{latex'}
\tikzset{->/.style = {decoration={markings,
                                  mark=at position 0.999
                                  with {\arrow[scale=2]{latex'}}},
                      postaction={decorate}}}
\tikzset{<-/.style = {decoration={markings,
                                  mark=at position 0.001 with {\arrowreversed[scale=2]{latex'}}},
                      postaction={decorate}}}
\tikzset{-->/.style = {decoration={markings,
                                  mark=at position 0.999
                                  with {\arrow[scale=2]{latex'[sep=0pt -2.5]}}},
                      postaction={decorate}}}
\tikzset{<--/.style = {decoration={markings,
                                  mark=at position 0.001
                                  with {\arrowreversed[scale=2]{latex'[sep=0pt -2.5]}}},
                      postaction={decorate}}}
\tikzset{<->/.style = {decoration={markings,
                                  mark=at position 0.001 with {\arrowreversed[scale=2]{latex'}},
                                  mark=at position 0.999 with {\arrow[scale=2]{latex'}},},
                      postaction={decorate}}}
\tikzset{->-/.style = {decoration={markings,
    mark=at position 0.50+0.225/#1 with {\arrow[scale=2]{latex'}}},
                      postaction={decorate}}}
\tikzset{-<-/.style = {decoration={markings,
                                   mark=at position 0.50-0.225/#1 with {\arrowreversed[scale=2]{latex'}}},
                       postaction={decorate}}}
\tikzset{->>/.style = {decoration={markings,
                                  mark=at position 1 with {\arrow[scale=2]{latex'}}},
                      postaction={decorate}}}
\tikzset{<<-/.style = {decoration={markings,
                                  mark=at position 0 with {\arrowreversed[scale=2]{twolatex'}}},
                      postaction={decorate}}}
\tikzset{<<->>/.style = {decoration={markings,
                                   mark=at position 0 with {\arrowreversed[scale=2]{twolatex'}},
                                   mark=at position 1 with {\arrow[scale=2]{twolatex'}}},
                       postaction={decorate}}}
\tikzset{->>-/.style = {decoration={markings,
                                   mark=at position 0.50+0.450/#1 with {\arrow[scale=2]{twolatex'}}},
                       postaction={decorate}}}
\tikzset{-<<-/.style = {decoration={markings,
                                   mark=at position 0.50-0.450/#1 with {\arrowreversed[scale=2]{twolatex'}}},
                       postaction={decorate}}}

\pgfarrowsdeclare{biggertip}{biggertip}{%
  \setlength{\arrowsize}{1pt}
  \addtolength{\arrowsize}{.1\pgflinewidth}
  \pgfarrowsrightextend{0}
  \pgfarrowsleftextend{-5\arrowsize}
}{%
  \setlength{\arrowsize}{1pt}
  \addtolength{\arrowsize}{.1\pgflinewidth}
  \pgfpathmoveto{\pgfpoint{-5\arrowsize}{4\arrowsize}}
  \pgfpathlineto{\pgfpointorigin}
  \pgfpathlineto{\pgfpoint{-5\arrowsize}{-4\arrowsize}}
  \pgfusepathqstroke
}
\tikzset{
	EdgeStyle/.style = {>=biggertip}
}

\tikzset{circ/.style = {fill, circle, inner sep = 0, minimum size = 3}}
\tikzset{scirc/.style = {fill, circle, inner sep = 0, minimum size = 1.5}}
\tikzset{mstate/.style={circle, draw, black, text=black, minimum width=0.7cm}}

\tikzset{eqpic/.style={baseline={([yshift=-.5ex]current bounding box.center)}}}

\definecolor{mblue}{rgb}{0.2, 0.3, 0.8}
\definecolor{morange}{rgb}{1, 0.5, 0}
\definecolor{mgreen}{rgb}{0, 0.4, 0.2}
\definecolor{mred}{rgb}{0.5, 0, 0}

% topology
\DeclareMathOperator{\dfr}{dfr} % deformation retract
\newcommand{\Cells}{\text{Cells}}
\newcommand{\RP}{\mathbb{RP}}
\newcommand{\CP}{\mathbb{CP}}

% order theory
\newcommand{\join}{\vee}
\newcommand{\meet}{\wedge}

% algebraic geometry
\DeclareMathOperator{\Rad}{Rad} % Radical of modules
\DeclareMathOperator{\Spec}{Spec}
% \renewcommand{\div}{\mathrm{div}} % divisors
\DeclareMathOperator{\Mer}{Mer} % Meromorphic functions
\DeclareMathOperator{\Exc}{Exc} % Exceptional divisor
\DeclareMathOperator{\Td}{Td} % Todd classes

% algebra

\DeclareMathOperator{\trdeg}{tr.deg.}
\DeclareMathOperator{\odd}{odd}
\DeclareMathOperator{\even}{even}
\DeclareMathOperator{\lcm}{lcm}
\DeclareMathOperator{\ord}{ord}
\DeclareMathOperator{\codim}{codim}
\DeclareMathOperator{\Ann}{Ann}
\DeclareMathOperator{\Ext}{Ext}
\DeclareMathOperator{\Sym}{Sym}
\DeclareMathOperator{\Out}{Out}
\DeclareMathOperator{\Aut}{Aut}
\DeclareMathOperator{\End}{End}
\DeclareMathOperator{\Inn}{Inn}
\DeclareMathOperator{\Mat}{Mat}
\DeclareMathOperator{\Fix}{Fix}
\DeclareMathOperator{\std}{std}
\DeclareMathOperator{\sgn}{sgn}
\DeclareMathOperator{\adj}{adj}
\DeclareMathOperator{\id}{id}
\DeclareMathOperator{\op}{op}
\DeclareMathOperator{\tr}{tr}
\DeclareMathOperator{\diag}{diag}
\DeclareMathOperator{\rank}{rank}
\DeclareMathOperator{\rk}{rk}
\DeclareMathOperator{\coker}{coker}
\DeclareMathOperator{\Mult}{Mult} % multiplier algebra
\DeclareMathOperator{\Frac}{Frac} % fraction field
\DeclareMathOperator{\Rat}{Rat}
\DeclareMathOperator{\Gal}{Gal}
\newcommand{\ioplus}{\overset{\text{int}}{\oplus}} % interior direct sum
\newcommand{\GG}{\mathcal G} % Galois correspondence
\newcommand{\FF}{\mathcal F} % Galois correspondence
\newcommand{\cupp}{\smile} % cup product
\newcommand{\capp}{\frown} % cap product
\newcommand{\actson}{\curvearrowright}
\newcommand{\bra}{\left\langle}
\newcommand{\ket}{\right\rangle}
\newcommand{\idealin}{\triangleleft\,}
\newcommand{\normalin}{\triangleleft\,}
\newcommand{\ip}[2]{\left\langle #1, #2 \right\rangle}
\newcommand{\bigslant}[2]
{{\raisebox{.2em}{$#1$}\left/\raisebox{-.2em}{$#2$}\right.}}
\newcommand{\properideal}{%
  \mathrel{\ooalign{$\lneq$\cr\raise.22ex\hbox{$\lhd$}\cr}}}

% categories
\newcommand{\cat}[1]{\mathbf{#1}}
\newcommand{\Ch}{\cat{Ch}} % Chain complexes   
\newcommand{\Set}{\cat{Set}}
\newcommand{\Grp}{\cat{Grp}}
\newcommand{\Ab}{\cat{Ab}}
\newcommand{\Ring}{\cat{Ring}}
\newcommand{\CRing}{\cat{CRing}}
\newcommand{\Top}{\cat{Top}}
\newcommand{\Vect}{\cat{Vect}}
\newcommand{\cmod}{\cat{mod}}

%% matrix groups
\newcommand{\GL}{\mathrm{GL}}
\newcommand{\Or}{\mathrm{O}}
\newcommand{\PGL}{\mathrm{PGL}}
\newcommand{\PSL}{\mathrm{PSL}}
\newcommand{\PSO}{\mathrm{PSO}}
\newcommand{\PSU}{\mathrm{PSU}}
\newcommand{\SL}{\mathrm{SL}}
\newcommand{\msl}{\mathrm{sl}}
\newcommand{\SO}{\mathrm{SO}}
\newcommand{\Spin}{\mathrm{Spin}}
\newcommand{\Sp}{\mathrm{Sp}}
\newcommand{\SU}{\mathrm{SU}}
\newcommand{\U}{\mathrm{U}}
\let\char\relax
\newcommand{\char}{\text{char}}

% analysis
\renewcommand{\Re}{\mathrm{Re}}
\renewcommand{\Im}{\mathrm{Im}}
\NewDocumentCommand{\dx}{o}{
  \IfNoValueTF{#1}
    {\dif x}                  % if no optional argument
    {\dif^{#1} \!x}         % if optional argument
}
\NewDocumentCommand{\dy}{o}{
  \IfNoValueTF{#1}
    {\dif y}                  % if no optional argument
    {\dif^{#1} \!y}         % if optional argument
}
\newcommand{\dt}{\dif t}
\newcommand{\du}{\dif u}
\newcommand{\dv}{\dif v}
\newcommand{\dz}{\dif z}
\newcommand{\ds}{\dif s}
\newcommand{\dw}{\dif w}
\newcommand{\dr}{\dif r}
\newcommand{\dm}{\dif m}
\newcommand{\df}{\dif f}
\newcommand{\dV}{\dif V}
\newcommand{\dS}{\dif S}
\newcommand{\dA}{\dif \mathrm{A}}
\newcommand{\dmu}{\dif \mu}
\newcommand{\dnu}{\dif \nu}
\newcommand{\dtheta}{\dif \theta}
\newcommand{\domega}{\dif \omega}
\newcommand{\dsigma}{\dif \sigma}
\newcommand{\dtau}{\dif \tau}
\newcommand{\dvarphi}{\dif \varphi}
\renewcommand{\dif}{\mathrm{d}}
\renewcommand{\Dif}{\mathrm{D}}
\renewcommand{\triangle}{\Delta}
\newcommand{\ptriangle}{\partial \triangle}
\DeclareMathOperator{\im}{im}
\DeclareMathOperator{\cis}{cis}
\DeclareMathOperator{\rot}{rot}
\DeclareMathOperator{\vspan}{span}
\DeclareMathOperator{\grad}{grad}
\DeclareMathOperator{\curl}{curl}
\DeclareMathOperator{\esssup}{esssup}
\newcommand{\hodge}{{\mathnormal{\star}}}
\DeclareMathOperator{\range}{range}
\DeclareMathOperator{\Hom}{Hom}
\DeclareMathOperator{\Int}{Int}
\DeclareMathOperator{\Conv}{Conv} % convex hull
\newcommand{\ind}{\mathbbm{1}}
\DeclareMathOperator{\Res}{Res}
\DeclareMathOperator{\graph}{graph}
\DeclareMathOperator{\diam}{diam}
\DeclareMathOperator{\supp}{supp}
\DeclareMathOperator{\Vol}{Vol} % Volume
\DeclareMathOperator{\Area}{Area} % Area
\DeclareMathOperator{\area}{area}
\DeclareMathOperator{\Ind}{Ind} % Index
\newcommand{\length}{\mathrm{length}}

% logic
\DeclareMathOperator{\MOD}{MOD}
\DeclareMathOperator{\Theory}{Theory}
\newcommand{\notimplies}{%
  \mathrel{{\ooalign{\hidewidth$\not\phantom{=}$\hidewidth\cr$\implies$}}}}

% nice
\newcommand{\half}{\frac{1}{2}}
\newcommand{\pair}{\del}
\newcommand{\taking}[1]{\xrightarrow{#1}}
\newcommand{\otaking}[1]{\xleftarrow{#1}}
\newcommand{\inv}{^{-1}}
\newcommand{\ot}{\leftarrow}
\newcommand{\ninfty}{-\infty}
\newcommand{\pinfty}{+\infty}
\newcommand{\floor}[1]{\left\lfloor #1 \right\rfloor}
\newcommand{\ceil}[1]{\left\lceil #1 \right\rceil}
\newcommand{\viff}{\Big\Updownarrow} % vertical iff
\newcommand{\bmid}{\,\middle\vert\,} % big mid

% probability
\newcommand{\Prob}{\mathbf{P}}
\renewcommand{\vec}[1]{\boldsymbol{\mathbf{#1}}}
\DeclareMathOperator{\Bin}{Bin}
\DeclareMathOperator{\Geo}{Geo}
\DeclareMathOperator{\Poi}{Poi}
\DeclareMathOperator{\Exp}{Exp}
\DeclareMathOperator{\Var}{Var} % Variance
\DeclareMathOperator{\Cov}{Cov}

% special letters
\newcommand{\N}{\mathbb{N}}
\newcommand{\Z}{\mathbb{Z}}
\newcommand{\Q}{\mathbb{Q}}
\newcommand{\R}{\mathbb{R}}
\newcommand{\C}{\mathbb{C}}
\newcommand{\F}{\mathbb{F}}
\newcommand{\E}{\mathbf{E}}
\newcommand{\D}{\mathbb{D}}
\renewcommand{\H}{\mathbb{H}}
\newcommand{\ps}{\mathcal{P}}
\newcommand{\M}{\mathcal{M}}
\renewcommand{\L}{\mathcal{L}}
\renewcommand{\O}{\mathcal{O}}
\renewcommand{\U}{\mathcal{U}}
\newcommand{\Omicron}{O}
\newcommand{\powerset}{\mathcal{P}}

% text
\newcommand{\st}{\text{ s.t. }}
\newcommand{\tand}{\quad \text{and} \quad}
\newcommand{\tor}{\quad \text{or} \quad}
\newcommand{\stand}{\text{ and }}
\newcommand{\stor}{\text{ or }}
\renewcommand{\tt}[1]{\textnormal{\textbf{(#1).}}} %tt=theorem title GET RID OF

% title format
\title{\textbf{\ncourse}}

\ifx \ntype \notes
  \author{Notes taken by \nauthor\, \nemail \\ Based on lectures by \nlecturer}
  \date{\nterm\ \nyear}
\else
  \ifx \ntype \problems
    \author{These problems were solved by \nauthor \\ \nemail}
  \else
  \fi
\fi

\makeatother


\begin{document}
\maketitle

% Insert cool image here

\newpage
\tableofcontents
\newpage

\section{Introduction to Probability}

\begin{exercise}
  Let $A$ and $B$ be independent events.
  Prove that $A$ and $B^c$
  are independent and that $A^c$ and $B^c$ are independent.
\end{exercise}
\begin{solution}
  By definition of indepence we have that
  \[
    \Prob(A \cap B) = \Prob(A) \Prob(B).
  \]
  Now, by the basic properties of a probabiliy function we get that:
  \[
    \Prob\del{A \cap B^c} =
    \Prob\del{A \cap (\Omega \setminus B)} =
    \Prob\del{A \setminus B} =
    \Prob\del{A \setminus (A \cap B)}.
  \]
  Since $A \cap B \subseteq A$ we get by independency of $A$ and $B$:
  \begin{align*}
    \Prob\del{A \cap B^c} &=
    \Prob\del{A \setminus (A \cap B)} =
    \Prob\del{A} - \Prob\del{A \cap B} =
    \Prob\del{A} - \Prob(A) \Prob(B) \\ &=
    \Prob\del{A} \del{1 - \Prob(B)} =
    \Prob\del{A} \Prob(B^c),
  \end{align*}
  so $A$ and $B^c$ are independent.
  Similarly, we have that $A^c$ and $B$ are independent.
  
  Next, by using the previous part of the exercise,
  and the fact that:
  \[
    A^c \cap B^c =
    \Omega \setminus \del{(A^c \cap B) \sqcup (A \cap B^c) \sqcup (A \cap B)}
  \]
  we get that:
  \begin{align*}
    \Prob\del{A^c \cap B^c} &=
    1 - \del{\Prob\del{A^c \cap B} + \Prob\del{A \cap B^c} + \Prob\del{A \cap B}} \\ &=
    1 - \del{\Prob(A) \Prob(B^c) + \Prob(A^c) \Prob(B) + \Prob(A) \Prob(B)} \\ &=
    1 - \del{(1 - \Prob(A^c)) \Prob(B^c) + \Prob(A^c) (1 - \Prob(B^c) +
      (1 - \Prob(A^c) (1 - \Prob(B^c)} \\ &=
    1 - \del{1 - \Prob(A^c) \Prob(B^c)} \\ &=
    \Prob(A^c) \Prob(B^c),
  \end{align*}
  which completes the proof.
\end{solution}

\begin{exercise}
  Suppose there are $r$ boxes and $r + 1$ balls.
  Balls are placed in the boxes at random, one at a time until,
  for the first time, one box has two balls.
  Find the probability that this takes place with the $n$th ball,
  for $2 \le n \le r + 1$.
\end{exercise}
\begin{solution}
  If the first collision was on the $n$th ball,
  this means that we had $n-1$ consecutive ball insertions
  without collisions, and then a collision.
  Note that the probability of a non-collision
  given that we inserted $k$ balls without collisions is
  equal to $\frac{r - k}{r}$, since exactly $k$ out of $r$ boxes
  are occupied.
  This implies that:
  \[
    \Prob(n\text{th ball collision}) =
    \Prob(\text{collision} \mid n-1\text{ non-collisions}) \cdot
    \prod_{k=0}^{n-2} \frac{r - k}{r},
  \]
  and since the probability of a collision given $n-1$ previous
  non-collision is exactly $\frac{n-1}{r}$,
  we conclude that
  \[
    \Prob(n\text{th ball collision}) =
    \frac{n - 1}{r} \cdot
    \prod_{k=0}^{n-2} \frac{r - k}{r}.
  \]
\end{solution}

\begin{exercise}
  A box has $10$ red balls and $5$ black balls.
  A ball is selected from the box.
  If the ball is red, it is returned to the box.
  If the ball is black,
  it is returned to the box along with two additional black balls.
  After this, another ball is drawn.
  \begin{enumerate}
    \item[(1)] What is the probability that this other ball is red?
    \item[(2)] If this other ball is black, what is the probability that
      the first ball selected was also black?
  \end{enumerate}
\end{exercise}
\begin{solution}\phantom{}
  \begin{enumerate}
    \item[(1)] Denote the event that the first ball drawn is red
      by $R_1$, the event its black by $B_1$,
      and similarly for the second ball we denote $R_2$ and $B_2$
      respectively.
      By using law of total probability we get that:
      \[
        \Prob(R_2) =
        \Prob(R_1) \Prob(R_2 \mid R_1) + \Prob(B_1) \Prob(R_2 \mid B_1) =
        \frac{10}{15} \cdot \frac{10}{15} +
        \frac{5}{15} \cdot \frac{12}{17} =
        \frac{4}{9} + \frac{10}{51} =
        \frac{98}{153}.
      \]
    \item[(2)] We are asked to calculate the probability of $\Prob(B_1 \mid B_2)$.
      By Bayes' theorem we get that:
      \[
        \Prob(B_1 \mid B_2) =
        \frac{\Prob(B_2 \mid B_1) \Prob(B_1)}{\Prob(B_2)},
      \]
      and by calculations from part $(1)$ we conclude:
      \[
        \Prob(B_1 \mid B_2) =
        \frac{\Prob(B_2 \mid B_1) \Prob(B_1)}{\Prob(B_2)} =
        \frac{\frac{7}{17} \cdot \frac{5}{15}}{1 - \frac{98}{153}} =
        \frac{21}{55}.
      \]
  \end{enumerate}
\end{solution}

\begin{exercise}
  In the game of bridge there are four players $A$, $B$, $C$ and $D$.
  Players $A$ and $C$ are partners and players $B$ and $D$ are partners.
  Each player gets $13$ cards.
  If players $A$ and $C$ have $9$ spades between them,
  what is the probability that the $4$ other spades are split three
  and one between the two other players?
\end{exercise}
\begin{solution}
  We want to find the probability of:
  \[
    \Prob(A \stand C \text{ have $4$ spades split $3:1$} \mid
      \text{$B$ and $D$ have $9$ spades}),
  \]
  but note that this reduces to the probability of splitting $26$
  cards with exactly $4$ spades in them between two players with the spades
  splitting exactly $3 : 1$ between the players.
  Also note that the hand of $A$ is detemined completely by the hand of $C$
  (similarly the hand of $C$ is detemined by the hand of $C$).
  Since we are choosing $13$ cards out of $26$ cards that contain $9$
  spades, the total number of possible hands for $A$ is $\binom{26}{13}$,
  and the number of hands containing exactly $3$ spades is
  $\binom{4}{3} \cdot \binom{22}{10}$.
  Since the probability of each hand is equally likely,
  we get that $\Prob$ is the uniform probability function,
  and thus:
  \[
    \Prob(A \stand C \text{ have $4$ spades split $3:1$}) =
    \frac{\binom{4}{3} \cdot \binom{22}{10}}{\binom{26}{13}} +
    \frac{\binom{4}{3} \cdot \binom{22}{10}}{\binom{26}{13}} =
    \frac{286}{575},
  \]
  where we summed the same probability twince
  since we need to consider both the case
  where $A$ has $3$ spades and $C$ has $1$, and the symmetric case where
  $C$ has $3$ spades and $A$ has $1$ (we are allowed to sum the probability
  of these events since they are disjoint as it is impossible for $A$
  for example to have both $1$ spade and $3$ spades in the same hand).
\end{solution}

\begin{exercise}
  An elevator stops at $10$ floors and has $7$ people on it.
  We label the various possible outcomes as follows:
  if $6$ people get out on some particular floor and one person gets out
  on some other particular floor, we write the outcome as $(6,1)$.
  If $3$ people get out on one particular floor, two people get
  out on some other particular floor and two people get out on yet another
  particular floor, then we write the outcome as $(3,2,2)$, etc.
  Assume that each person behaves independently of the others and that
  each person has an equal chance of exiting on any floor.
  Calculate the probability of $(6,1)$ and the probability of $(3,2,2)$.
\end{exercise}
\begin{solution}
  NOTE: In this question we first assume order matters -
  For example, the events $(6,1)$ and $(1,6)$ are different,
  then show the solution in case it does not matter.

  The multinomial formula states the the number of options of assigning
  $7$ people to $10$ floors is exactly:
  \[
    \frac{7!}{a_1 ! a_2 ! \cdots a_{10}!},
  \]
  where $a_i$ denotes the amount of people who got off at the $i$th floor.
  This implies that for an event of the form $(a_1,\dots,a_{10})$,
  under the assumption from the exercise $\Prob$ is uniform
  and the events of people getting off are independent we have:
  \[
    \Prob(a_1,\dots,a_{10}) =
    \frac{7!}{a_1 ! a_2 ! \cdots a_{10}!} \cdot \frac{1}{10^7},
  \]
  where $10^{-7}$ denotes the total number of options,
  and in order to compute the probability of an event of the form $(6,1)$,
  we multiply by a constant $C$ which is the number of tuples $(a_1,\dots,a_{10})$,
  that possibly represent the event (for example $(7,0,\dots,0)$ and $(0,7,\dots,0)$
  and so on all represent the event $(7)$, so the constant is $10$).
  It is clear that this constant is $\binom{10}{k}$ where $k$ is the number
  of floors where more than one person got off, since we care about the order
  they got off.
  It follows that:
  \[
    \Prob(6,1) =
    \binom{10}{2} \cdot \frac{7!}{6! 1!} \cdot \frac{1}{10^7} =
    \frac{63}{2000000}.
  \]
  and
  \[
    \Prob(3,2,2) =
    \binom{10}{3} \cdot \frac{7!}{3! 2! 2!} \cdot \frac{1}{10^7} =
    \frac{63}{25000}.
  \]
  Note that if we did not distinguish between the events $(6,1)$
  and $(1,6)$ (i.e. in the event $(6,1)$ it is possible $6$ people got off
  after the first person got off),
  then we need to multiply the probabilities we got by the
  number of permutations of the multisets $\set{6,1}$ and $\set{3,2,2}$
  respetively. Since it is easy to see that they are $2$ and $3$
  respectively, we get that in that case
  \[
    \Prob(6,1) =
    2 \cdot \binom{10}{2} \cdot \frac{7!}{6! 1!} \cdot \frac{1}{10^7} =
    \frac{63}{1000000}.
  \]
  and
  \[
    \Prob(3,2,2) =
    3 \cdot \binom{10}{3} \cdot \frac{7!}{3! 2! 2!} \cdot \frac{1}{10^7} =
    \frac{189}{25000}.
  \]
\end{solution}

\begin{exercise}
  \label{ex:one-six}
  Let $a,b \in (0,1)$.
  Players $A$ and $B$ take turns throwing a dart at a small target.
  Each time player $A$ throws, the probability that she hits the target is $a$.
  Each time player $B$ throws, the probability that he hits the target is $b$.
  Assume that the outcomes on different throws are all independent of one another.
  Player $A$ goes first.
  The first person to hit the target wins the game.
  Show that the probability that player $A$ wins the game
  is $\frac{a}{1 - (1 - a) (1 - b)}$.
\end{exercise}
\begin{solution}
  Let $W$ denote the event that player $A$ wins,
  and let $W_i$ be denote the event that player $A$ wins on her $i$th toss.
  Since the events $\set{W_i}_{i=1}^{\infty}$ are disjoint
  and $W = \bigcup_{i=1}^{\infty} W_i$ we get:
  \[
    \Prob(W) =
    \sum_{i=1}^{\infty} \Prob(W_i).
  \]
  It is clear that $W_i$ is exactly the event in which
  $A$ and $B$ miss their $1,\dots,i-1$ shots,
  and then $A$ hits the mark.
  Since the outcomes of all throws are independent of one another,
  we get that:
  \begin{align*}
    \Prob(W_i) &=
    \Prob(\text{$A$ $1$st shot miss})
    \Prob(\text{$B$ $1$st shot miss}) \cdots
    \Prob(\text{$B$ $i-1$th shot miss})
    \Prob(\text{$A$ $i$th shot hit}) \\ &=
    (1 - a)^{i-1} (1 - b)^{i-1} a =
    a \del{(1 - a)(1 - b)}^{i-1}.
  \end{align*}
  It follows that:
  \[
    \Prob(W) =
    a \sum_{i=1}^{\infty} (1 - a - b + ab)^{i-1}.
  \]
  This is an infinite geometric series with factor $(1 - a)(1 - b)$,
  and thus we conclude that:
  \[
    \Prob(W) =
    a \cdot \frac{1}{1 - (1 - a) (1 - b)} =
    \frac{a}{1 - (1 - a) (1 - b)},
  \]
  which concludes the proof.
  
  However, in order to solve \Cref{ex:one-seven},
  we additionally compute the probability for $A$ to win given that
  $B$ started the game.
  Now since $B$ has one additional shot at the start we get
  \[
    \Prob(W_i) =
    (1 - a)^{i-1} (1 - b)^{i} a =
    a (1 - b) \del{(1 - a)(1 - b)}^{i-1}.
  \]
  It is clear that now we will get
  $\Prob(W) = \frac{a (1 - b)}{1 - (1 - a) (1 - b)}$.
\end{solution}

\begin{exercise}
  \label{ex:one-seven}
  A coin is tossed repeatedly.
  On each toss, the probability of heads is $p$ and the probability of tails
  is $1 - p$, where $p \in (0,1)$.
  A head run of length $r$ is a sequence of $r$ heads in a row.
  For example, if one sees $HHTTHTHTTHHHTHHHHHHT\cdots$,
  then a head run of length $6$ occurred starting on the 14th toss.
  Calculate the probability that a head run of length $r$ will occur before
  a tail run of length $s$.
  Hint: Analyze the problem in such a way that you can use the formula
  derived in \Cref{ex:one-six}.
\end{exercise}
\begin{solution}
  Suppose $A$ and $B$ are players who flip the same coin from the original
  question, with probability $p \in (0,1)$ for heads, and probability
  $1 - p$ for tails.
  Before the game starts, a third unrelated player flips the coin once,
  and then if the result is heads,
  player $A$ starts flipping it until the coin shows tails,
  in which case he lets $B$ flip the coin until he flips heads,
  and so on and so forth.
  If the first coin showed tails, the same game is played only player $B$
  starts instead of player $A$.
  We say that player $A$ wins if he flipped $r - 1$ heads in a row,
  and that player $B$ wins if he flipped $s - 1$ tails in a row.
  Now the question becomes what is the probability that $A$ wins?
  The win condition of $A$ has probability $p^{r-1}$ and the win
  condition of $B$ is $(1-p)^{s-1}$.
  It is not hard to notice that the probability we are looking
  for in the original question is exactly the probability that player $A$
  wins the game (though it is relatively hard to think of this game
  in the first place).
  From the law of total probability we get:
  \[
    \Prob(A \text{ wins}) =
    \Prob(A \text{ wins} \mid \text{$1$st flip } H) \Prob(\text{$1$st flip } H) +
    \Prob(A \text{ wins} \mid \text{$1$st flip } T) \Prob(\text{$1$st flip } T).
  \]
  From \Cref{ex:one-six} we have that:
  \begin{align*}
    \Prob(A \text{ wins} \mid \text{$1$st flip }H) &=
    \frac{p^{r-1}}{1 - (1 - p^{r-1}) (1 - (1-p)^{s-1})} \\
    \Prob(A \text{ wins} \mid \text{$1$st flip }T) &=
    \frac{\del{1 - (1-p)^{s-1}} p^{r-1}}{1 - (1 - p^{r-1}) (1 - (1-p)^{s-1})}
  \end{align*}
  and thus:
  \begin{align*}
    \Prob(A \text{ wins}) &=
    \frac{p^{r}}{1 - (1 - p^{r-1}) (1 - (1-p)^{s-1})} +
    \frac{(1 - (1-p)^{s-1}) p^{r-1} (1 - p)}{1 - (1 - p^{r-1}) (1 - (1-p)^{s-1})} \\ &=
    \frac{p^{r - 1} \del{1 - (1-p)^s}}{1 - (1 - p^{r-1}) (1 - (1-p)^{s-1})}.
  \end{align*}
\end{solution}

\newpage

\section{Distribution and Expectation}

\begin{exercise}
  Let $X$ be uniformly distributed on $[n] = \set{1,\dots,n}$.
  That is, $\Prob(X = j) = \frac{1}{n}$,
  for all $j \in [j]$.
  Calculate the variance of $X$.
\end{exercise}
\begin{solution}
  We have seen that $\Var(X) = \E[X^2] - \E[X]^2$.
  By definition of expectation of a discrete random variable we have that
  \[
    \E[X] =
    \sum_{k=1}^{n} k \cdot \Prob_X(k) =
    \frac{1}{n} \sum_{k=1}^{n} k =
    \frac{n + 1}{2},
  \]
  and therefore $\E[X]^2 = (\frac{n + 1}{2})^2$
  and similarly by the formula for $\sum_{k=1}^{n} k^2$ we get:
  \[
    \E(X^2) :=
    \sum_{k=1}^{n} k^2 \cdot \Prob_{X^2}(k^2) =
    \frac{1}{n} \sum_{k=1}^{n} k^2 =
    \frac{(n + 1)(2n + 1)}{6},
  \]
  so finally:
  \[
    \Var(X) =
    \E[X^2] - \E[X]^2 =
    \frac{(n + 1)(2n + 1)}{6} - \del{\frac{n + 1}{2}}^2 =
    \frac{n^2 - 1}{12}.
  \]
\end{solution}

\begin{exercise}
  Solve the following:
  \begin{enumerate}
    \item[(1)] Prove that:
      \[
        \sum_{k=0}^{n} \binom{n}{k}^2 =
        \binom{2n}{n}.
      \]
    \item[(2)] Two people toss a fair coin $n$ times. 
      What is the probability that they get the same number of heads?
      % use (1)
    \item[(3)] Let $p_n$ denote the answer from $(2)$.
      Use Stirling's approximation in order to show that
      $\lim_{n \to \infty} \sqrt{n} p_n = \frac{1}{\sqrt{\pi}}$.
  \end{enumerate}
\end{exercise}
\begin{solution} \phantom{}
  \begin{enumerate}
    \item[(1)] On the right hand side of the equation
      is clearly the number of ways to choose $n$ elements
      out of $2n$ elements.
      On the left hand side we have that:
      \[
        \sum_{k=0}^{n} \binom{n}{k}^2 =
        \sum_{k=0}^{n} \binom{n}{k}\binom{n}{n - k},
      \]
      which is the sum of the amounts of options to pick $k$
      elements out of $n$ elements, and $n-k$ elements out of $n$ different
      elements, for all $0 \le k \le n$.
      Since these amount must be equal (each choice of $n$ elements
      out of $2n$ elements is some choice of $k$ elements out of $n$
      elements and $n-k$ elements out of $n$ different elements and
      vice versa), the equality is satisfied.
    \item[(2)] Let $X$ and $Y$ be random variables that denote
      the amount of heads each of the two people got after $n$ tosses
      respectively.
      We then have by part $(1)$ that:
      \[
        \Prob(X = Y) =
        \sum_{k=0}^{n} \Prob(X = k) \cdot \Prob(Y = k) =
        \del{\frac{1}{2^n}}^2 \sum_{k=0}^{n} \binom{n}{k}^2 =
        \frac{\binom{2n}{n}}{4^{n}}.
      \]
    \item[(3)] We need to show that:
      \[
        \lim_{n \to \infty} \frac{\binom{2n}{n}\sqrt{n}}{4^{n}} = \frac{1}{\sqrt{\pi}}.
      \]
      From Stirling's approximation for sufficiently large $n$ we get
      that:
      \[
        \frac{\binom{2n}{n}\sqrt{n}}{4^{n}} =
        \frac{\frac{(2n)!}{n!n!}\sqrt{n}}{4^{n}} \approx
        \frac{(2n)^{2n} e^{-2n} 2\sqrt{\pi n}}{n^{2n} e^{-2n}2\pi n}
        \frac{\sqrt{n}}{4^{n}} =
        \frac{(2n)^{2n} }{n^{2n} \sqrt{\pi} 4^{n}} =
        \frac{1}{\sqrt{\pi}},
      \]
      Indeed, we then get that:
      \[
        \lim_{n \to \infty} \sqrt{n} p_n =
        \frac{1}{\sqrt{\pi}},
      \]
      which completes the proof.
  \end{enumerate}
\end{solution}

\begin{exercise}
  In a sequence of Bernoulli trials the probability of success is $p$ and the
  probability of failure is $1 - p$.
  Write down an expression for the probability of getting $N$ successes
  before getting $M$ failures.
  % (Hint: Let Aj denote the
  % event that when the N -th success occurs, there are already exactly j failures.
  % If you can calculate P (Aj ), then you can easily write down an expression
  % for the probability of getting N successes before getting M failures. To
  % calculate P (Aj ), just consider what happens on the first j + N − 1 trials.)
\end{exercise}
\begin{solution}
  Let $q = 1 - p$, and let $A_j$ denote the event that when the $N$th success
  occurs, there are already exactly $j$ failures,
  for $0 \le j \le M-1$.
  Note that in the event $A_j$ we simply have to choose when the $j$
  failures would occur (they could occur anywhere from the $1$st to the
  $N + j - 1$th toss, since the $N + j$th toss is necessarily a win).
  Since in each such choice the probability of the event is exactly
  $p^{N} q^j$ we get that:
  \[
    \Prob(A_j) =
    \binom{N + j - 1}{j} p^N q^j.
  \]
  The event that $N$ successes occur before $M$ failures is
  exactly the disjoint union of all events of the form $A_j$
  so by the properties of the probability function we get:
  \[
    \Prob(\text{$N$ successes before $M$ failures}) =
    \Prob\del{\bigsqcup_{j=0}^{M-1} A_j} =
    \sum_{j=0}^{M-1} \binom{N + j - 1}{j} p^N q^j.
  \]
\end{solution}

\begin{exercise}
  \label{two-four}
  Let $A \subset \set{2, 3, 4, 5, 6}$.
  Consider the following game.
  You toss a die until you get either a $1$ or an element from $A$.
  Your score in the game is the value of your last toss.
  What set $A$ should you choose if you want to
  maximize the expected value of your score?
\end{exercise}
\begin{solution}
  Let $S = A \cup \set{1}$ denote the set of sides which stop the game,
  and let $|S| = m$.
  The probability that we stop in each roll is exactly
  $r := \frac{m}{6}$, and since all rolls are independent and identically
  distributed, if we denote the number of the terminal roll by $T$,
  we have $T \sim \Geo(r)$, and by the table of expectations of random
  variables we get $\E[T] = \frac{1}{r}$.
  Denote the value on the terminal roll by $X_T$ (this is our score).
  It takes values in $k \in S$ with equal probability (uniform on $S$),
  and therefore:
  \[
    \E[X_T] =
    \sum_{k \in S} k \cdot \Prob(X_T = k) =
    \frac{1}{m} \sum_{k \in S} k.
  \]
  We now need to find the set $A$ which would maximize this value.
  At first, let $A = \emptyset$.
  We see that whenever adding any number to $A$,
  the value of $\frac{1}{m}$ changes regardless of what number we added,
  and $\sum_{k \in S} k$ increases more as the number we choose to add is bigger.
  It follows that for a set of size $|A| = k$,
  the best $A$ we can choose is $A = \set{6,\dots,6-k+1}$.
  Thus we only need to find the set $A$ of these forms which maximizes $\E[X_T]$.
  We have that:
  \[
    \E[X_T] =
    \frac{1}{m} \sum_{k \in S} k =
    \begin{cases}
      1, &A = \emptyset, \\
      3.5, &A = \set{6}, \\
      4, &A = \set{6,5}, \\
      4, &A = \set{6,5,4}, \\
      \frac{19}{5}, &A = \set{6,5,4,3}, \\
      \frac{21}{6}, &A = \set{6,5,4,3,2}.
    \end{cases}
  \]
  A set $A = \set{5,6}$, or equivalently $A = \set{4,5,6}$
  maximizes our score.
  In both cases the expectation is $4$.
\end{solution}

\begin{exercise}
  Consider the game in \Cref{two-four} with the following modification:
  your score for the game is the value of your last toss minus the
  number of tosses.
  Now how should you choose $A$ in order to maximize the expected value of
  your score?
\end{exercise}
\begin{solution}
  In this game the expected score is:
  \[
    \E[\text{score}] =
    \E[X_T] - \E[T] =
    \frac{1}{m} \sum_{k \in S} k - \frac{6}{m} =
    \frac{\sum_{k \in S} k - 6}{m}
  \]
  where the last transition follows from calculations from \Cref{two-four}.
  Since $m = 1 + |A|$ and we always have $1 \in S$ we
  even get that:
  \[
    \E[\text{score}] =
    \frac{\sum_{k \in S} k - 6}{m} =
    \frac{\sum_{a \in A} a - 5}{1 + |A|}.
  \]
  From the same reasoning as in \Cref{two-four},
  we only need to consider the sets of the form $\set{6,\dots,6-|A|+1}$,
  and we compute to get:
  \[
    \E[\text{score}] =
    \frac{\sum_{a \in A} a - 5}{1 + |A|} =
    \begin{cases}
      -5, &A = \emptyset, \\
      0.5, &A = \set{6}, \\
      2, &A = \set{6,5}, \\
      2.5, &A = \set{6,5,4}, \\
      2.6, &A = \set{6,5,4,3}, \\
      2.5, &A = \set{6,5,4,3,2}.
    \end{cases}
  \]
  It follows that the maximum expected value is $2.6$ achieved by
  the set $A = \set{3,4,5,6}$.
\end{solution}

\begin{exercise}
  In class we have shown that if $X_{N,m,n}$ is distributed according to the
  hypergeometric distribution with parameters $N$, $m$, $n$,
  then $\Prob(X_{N,m,n} = j) = \frac{\binom{m}{j} \binom{N-m}{n-j}}{\binom{N}{n}}$
  for $j$ satisfying $0 \vee (n - N + m) \le j \le n \wedge m$.
  We explained intuitively that if $n$ is fixed and $N$ and $m$ are very large
  with $\frac{m}{N} \approx p$, then the distribution is close to $\Bin(n,p)$.
  Prove that for $p \in (0,1)$ we have:
  \[
    \lim_{N \to \infty,m \to \infty, \frac{m}{N} \to p} \Prob(X_{N,m,n} = j) =
    \lim_{N \to \infty,m \to \infty, \frac{m}{N} \to p} \frac{\binom{m}{j} \binom{N-m}{n-j}}{\binom{N}{n}} =
    \binom{n}{j} p^j (1 - p)^{n - j}.
  \]
  Hint: Let $N \to \infty$, let $m = m_N$ and define $p_N = \frac{m_N}{N}$.
  The assumption is that $\lim_{N \to \infty} p_N = p$.
\end{exercise}
\begin{solution}
  In this exercise we will write factorials using the falling–factorial notation
  suggested in wikiepdia $(a)_k = a(a-1)\cdots(a-k+1)$.
  We have
  \[
    \binom{m}{j}\binom{N-m}{n-j}
    = \frac{(m)_j}{j!}\cdot\frac{(N-m)_{n-j}}{(n-j)!}
    \tand
    \binom{N}{n} = \frac{(N)_n}{n!}.
  \]
  Thus we can write:
  \[
    \Prob\del{X_{N,m,n}=j}
    = \binom{n}{j}\,
      \frac{(m)_j (N-m)_{n-j}}{(N)_n}.
  \]
  We can divide the numerator and denominator by $N^n$ to get:
  \[
    \Prob\del{X_{N,m,n}=j}
    = \binom{n}{j}\,
      \frac{
        \del{(m)_j / N^j}\,\del{(N-m)_{n-j}/N^{\,n-j}}
      }{
        (N)_n / N^n
      }.
  \]
  For fixed $n$ and $j$ we get:
  \[
    \frac{(m)_j}{N^j}
    = \prod_{i=0}^{j-1} \frac{m-i}{N}
    = \prod_{i=0}^{j-1} \del{ p_N - \frac{i}{N} },
  \]
  where $p_N = \frac{m_N}{N}$.
  Since $p_N \to p$ and $\frac{i}{N} \to 0$,
  each factor tends to $p$, and therefore
  \[
    \lim_{N\to\infty} \frac{(m)_j}{N^j} = p^j.
  \]
  Similarly,
  \[
    \frac{(N-m)_{n-j}}{N^{\,n-j}}
    = \prod_{i=0}^{n-j-1} \del{ 1 - p_N - \frac{i}{N} }
    \taking{p_N \to p} (1-p)^{\,n-j}.
  \]
  Finally,
  \[
    \frac{(N)_n}{N^n}
    = \prod_{i=0}^{n-1} \del{ 1 - \frac{i}{N} }
    \taking{N \to \infty} 1.
  \]
  Putting everything together gives:
  \[
    \lim_{N \to \infty,m \to \infty, \frac{m}{N} \to p} \Prob(X_{N,m,n} = j) =
    \binom{n}{j} p^j (1 - p)^{n - j},
  \]
  which completes the proof.
\end{solution}

\newpage

\section{More Random Variables}

\begin{exercise}
  Consider an equilateral triangle with sides of length $s$.
  Choose a point randomly from one of the sides of the triangle.
  Let $X$ denote the distance between the point and the opposite vertex.
  Find the distribution function and the density for this random variable.
\end{exercise}
\begin{solution}
  We may fix a side of the triangle and embed it in $\R^2$ such that
  the side from which we sample the point is on the $x$-axis
  interval $I := [-s/2,s/2]$.
  The height of the triangle with respect to the side we chose (although
  since it is an equilateral triangle the side we choose does not matter)
  is $h = \frac{\sqrt 3}{2} s$ the Pythagorean theorem.
  Note that the opposite vertex from any point on our base side
  is $(0,h) \in \R^2$, and so if $T \sim U(I)$ is a random variable representing
  the $x$-coordinate of our point then we have the relation:
  \[
    X = \sqrt{T^2 + h^2},
  \]
  so $X$ has support $[h,s]$.
  In order to compute the distribution function we see that:
  \[
    F_X(x) =
    \Prob(X \le x) =
    \Prob(T^2 + h^2 \le x^2) =
    \Prob(|T| \le \sqrt{x^2 - h^2}),
  \]
  and since $T \sim U(I)$ we get
  \[
    F_X(x) =
    \Prob(|T| \le \sqrt{x^2 - h^2}) =
    \frac{2 \sqrt{x^2 - h^2}}{s} \quad \text{when } h \le x \le s.
  \]
  Then we get that:
  \[
    F_X(x) =
    \begin{cases}
      0 &x < h, \\
      \frac{2 \sqrt{x^2 - h^2}}{s} &h \le x \le s, \\
      1 &x > s.
    \end{cases}
  \]
  In order to compute the density function we just need to differentiate
  the function as follows:
  \[
    f_X(x) =
    \frac{\dif}{\dx} F_X(x) =
    \frac{2}{s} \cdot \frac{x}{\sqrt{x^2 - h^2}} \quad
      \text{when } h \le x \le s.
  \]
  Thus,
  \[
    f_X(x) =
    \begin{cases}
      \frac{2x}{s \sqrt{x^2 - h^2}} &h \le x \le s, \\
      0 &\text{otherwise}.
    \end{cases}
  \]
\end{solution}

\begin{exercise}
  Let $X$ be uniformly distributed on $[0,1]$.
  Let $Y = X^{\frac{1}{\beta}}$, where $\beta \neq 0$.
  Find the density function for $Y$.
\end{exercise}
\begin{solution}
  Since the map $g \colon x \mapsto x^{1/\beta}$ is continuous,
  we get that $Y = g(X) = X^{1/\beta}$ is a well defined random variable.
  Whenever $\beta > 0$ we get that the support of $Y$ is $[0,1]$ and
  that the is increasing, so for $0 \le y \le 1$ we have:
  \[
    F_Y(y) =
    \Prob(Y \le y) =
    \Prob(X^{1/\beta} \le y) =
    \Prob(X \le y^{\beta}) = y^{\beta},
  \]
  and by differentiating we get $f_Y(y) = F_Y'(y) = \beta y^{\beta - 1}$
  for $0 < y < 1$, we conclude that:
  \[
    f_Y(y) =
    \begin{cases}
      \beta y^{\beta - 1} &0 < y < 1, \\
      0 &\text{otherwise}.
    \end{cases}
  \]
  When $\beta < 0$ we can write $Y = X^{-1/\alpha}$ for $\alpha = -\beta > 0$.
  The support of $Y$ is now $[1,\infty)$ and the function is decreasing,
  therefore for $y \geq 1$,
  \[
    F_Y(y) =
    \Prob(Y \le y) =
    \Prob(X^{-1/\alpha} \le y) =
    \Prob(X \geq y^{-\alpha}) = 1 - y^{-\alpha}.
  \]
  Thus, similarly $f_Y(y) = F_Y'(y) = (-\beta) y^{\beta - 1}$ for $y > 1$,
  and in conclusion:
  \[
    f_Y(y) =
    \begin{cases}
      (-\beta) y^{\beta - 1} &y > 1, \\
      0 &\text{otherwise}.
    \end{cases}
  \]
\end{solution}

\begin{exercise}
  A random variable $X$ is know to be distributed according to a normal distribution.
  It is also known that $\Prob(X \le 160) = \half$ and $\Prob(X \le 140) = \frac{1}{4}$.
  Find $\mu$ and $\sigma$, and find $\Prob(X \geq 200)$.
\end{exercise}
\begin{solution}
  We know that $X \sim N(\mu,\sigma^2)$ for some values which we would like to find.
  Since $X$ distributes normally, we get that $\Prob(X \le \mu) = \half$,
  but we are given $\Prob(X \le 160)$ and thus $\mu = 160$.
  Then we can notice that if $Z \sim \N(0,1)$,
  then from the data given we can deduce that:
  \[
    \Prob\del{Z \le \frac{140 - \mu}{\sigma}} =
    \Prob\del{Z \le \frac{140 - 160}{\sigma}}
    \le \frac{1}{4},
  \]
  and since we can check with tables that the $25$th percentile of the
  standard normal distribution is approximately $-0.674$, we get that:
  \[
    \frac{-20}{\sigma} = -0.674 \implies
    \sigma = \frac{20}{0.674} \approx 29.7
  \]
  Finally, in order to compute $\Prob(X \geq 200)$ we just have to standardize
  again using $Z \sim N(0,1)$ as follows:
  \[
    \Prob(X \geq 200) =
    \Prob\del{Z \le \frac{200 - 160}{29.7}} =
    \Prob(Z \geq 1.35) \approx 0.088.
  \]
\end{solution}

\begin{exercise}
  Let $X$ be a continuous random variable with distribution function $F$
  and continuous density function $f$.
  Assume that $X$ is nonnegative; that is, $F(x) = 0$, for $x \le 0$.
  \begin{enumerate}
    \item[(1)] Assume that $\E[X] < \infty$.
      Prove that $\E[X] = \int_{0}^{\infty} (1 - F(x)) \dx$.
      (Hint: By assumption, $\int_{0}^{\infty} x f(x) \dx < \infty$,
      from which follows that $\lim_{x \to \infty} \int_{x}^{\infty} y f(y) \dy = 0$.
      Also, note that $1 - F(x) = \int_{x}^{\infty} f(y) \dy$.)
    \item[(2)] Assume that $\E[X] = \infty$.
      Prove that $\int_{0}^{\infty} (1 - F(x)) \dx = \infty$.
  \end{enumerate}
\end{exercise}
\begin{solution} \phantom{}
  \begin{enumerate}
    \item[(1)] We recall that the definition of expectation in this
      case is:
      \[
        \E[X] := \int_{0}^{\infty} x f(x) \dx < \infty.
      \]
      As suggested in the hint, we write:
      \[
        1 - F(x) = \int_{x}^{\infty} f(y) \dy.
      \]
      Using this identity we can write:
      \[
        \int_{0}^{\infty} \del{1 - F(x)} \dx =
        \int_{0}^{\infty} \del{\int_{x}^{\infty} f(y) \dy} \dx.
      \]
      Since the integrand is nonnegative we can change the order of integration
      (Fubini) and get:
      \[
        \int_{0}^{\infty} \del{\int_{x}^{\infty} f(y) \dy} \dx =
        \int_{x}^{\infty} f(y) \del{\int_{0}^{y} \dx} \dy =
        \int_{x}^{\infty} y f(y) \dy = \E[X],
      \]
      where the upper limit of the integral changed to $y$ since we
      are integrating over the domain
      $D = \set{(x,y) \in \R^2 \mid 0 \le x \le y < \infty}$.
      This completes the first part of the proof.
    \item[(2)] In case $\E[X] = \infty$, then:
      \[
        \E[X] := \int_{0}^{\infty} y f(y) \dy = \infty.
      \]
      We cannot use exactly the same argument from part $(1)$,
      but for every $M > 0$ we can do something similar and we get:
      \[
        \int_{0}^{M} \del{1 - F(x)} \dx =
        \int_{0}^{M} \del{\int_{x}^{\infty} f(y) \dy} \dx,
      \]
      where now the domain of integration is
      $D = \set{(x,y) \in \R^2 \mid 0 \le x M \stand x \le y < \infty}$.
      By switching the order of integration (by Fubini) we get:
      \begin{align*}
        \int_{0}^{M} \del{1 - F(x)} \dx &=
        \int_{0}^{M} \del{\int_{0}^{y} f(y) \dx} \dy + 
        \int_{M}^{\infty} \del{\int_{0}^{M} f(y) \dx} \dy \\ &=
        \int_{0}^{M} y f(y) \dy + 
        M \int_{M}^{\infty} f(y) \dy,
      \end{align*}
      where in the first step we decomposed the integration over the ``triangle''
      and the ``infinite at one side rectangle'' of the integration domain $D$.
      Since both terms on the right are nonnegative we have that:
      \[
        \int_{0}^{M} \del{1 - F(x)} \dx \geq
        \int_{0}^{M} y f(y) \dy,
      \]
      and this holds for any $M > 0$.
      Of course, as $M \to \infty$ we get by the comparison theorem that:
      \[
        \infty = \lim_{M \to \infty} \int_{0}^{M} y f(y) \dy \le
        \lim_{M \to \infty} \int_{0}^{M} \del{1 - F(x)} =
        \int_{0}^{\infty} (1 - F(x)) \dx,
      \]
      thus $\int_{0}^{\infty} (1 - F(x)) \dx = \infty$ which completes
      the proof.
  \end{enumerate}
\end{solution}

\begin{exercise}
  Let $X$ be a continuous random variable with density $f_X(x)$,
  and let $\psi \colon \R \to \R$ be a measurable function.
  Define $Y = \psi(X)$.
  In class, we saw that the expected value of $Y$ is
  $\E[Y] = \int_{-\infty}^{\infty} \psi(x) f_X(x) \dx$.
  Assume now that $f_X$ is continuous and vanishes off the interval $[0,1]$.
  Also assume now that $\psi$ is strictly increasing and continuously
  differentiable.
  In class we showed that $Y$ has a density function
  $f_Y(y) = \frac{f_X(x)}{\psi'(x)} = \frac{f_X(\psi^{-1}(y))}{\psi'(\psi^{-1}(y))}$,
  where $y = \psi(x)$.
  But then, by the definition of the expected value for a continuous
  random variable, we have $\E[Y] = \int_{-\infty}^{\infty} y f_Y(y) \dy$.
  Use a change of variables to show that these two definitions are consistent.
  That is, show that
  $\int_{-\infty}^{\infty} \psi(x) f_X(x) \dx = \int_{-\infty}^{\infty} y f_Y(y) \dy$.
\end{exercise}
\begin{solution}
  We know that $\psi$ is strictly increasing, continuously differentiable,
  and thus $\psi'$ is positive, and we can also perform a change of variables
  on the integral:
  \[
    \int_{-\infty}^{\infty} y f_Y(y) \dy =
    \int_{-\infty}^{\infty} y \frac{f_X(\psi^{-1}(y))}{\psi'(\psi^{-1}(y))} \dy.
  \]
  The change necessary in this problem is $y = \psi(x)$.
  Since $\psi(x)$ is injective and continuously differentiable we get:
  \[
    y = \psi(x) \tand
    \dy = \psi'(x) \dx \tand
    x = \psi^{-1}(y).
  \]
  Subtituting these into the integral we get that:
  \[
    \int_{-\infty}^{\infty} y f_Y(y) \dy =
    \int_{-\infty}^{\infty} y \frac{f_X(\psi^{-1}(y))}{\psi'(\psi^{-1}(y))} \dy =
    \int_{-\infty}^{\infty} \psi(x) \frac{f_X(x)}{\psi'(x)} \psi'(x) \dx,
  \]
  the Jacobian $\psi'(x)$ (which is always positive so we disregard absolute
  value) cancels and we conclude:
  \[
    \int_{-\infty}^{\infty} y f_Y(y) \dy =
    \int_{-\infty}^{\infty} \psi(x) \frac{f_X(x)}{\psi'(x)} \psi'(x) \dx =
    \int_{-\infty}^{\infty} \psi(x) f_X(x) \dx,
  \]
  which completes the proof.
\end{solution}

\begin{exercise}
  In real life we can't construct a continuous random variable, so we need
  to approximate it by a discrete one which takes on only a finite number
  of possible values.
  Let $X$ be uniformly distributed on $[0,1]$.
  Let $X_n$ be uniformly distributed on the $n$ points
  $\{\frac{j}{n}\}_{j=1}^{n}$.
  Let $H(x)$ be a continuous function from $[0,1]$ to $\R$.
  Explain why $\lim_{n \to \infty} \E[H(X_n)] = \E[H(X)]$.
\end{exercise}
\begin{solution}
  By definition of expectation of a discrete random variable we see that:
  \[
    \E\left[H(X_n)\right] = \frac{1}{n} \sum_{j=1}^{n} H\del{\frac{j}{n}},
  \]
  and on the other hand by the law of the unconscious statistician,
  \[
    \E[H(X)] =
    \int_{0}^{1} H(x) \dx.
  \]
  Note that since $H$ is a continuous function on $[0,1]$ to $\R$,
  so in particular, it is integrable,
  and its Riemann integral is the limit of the Riemann sums,
  and since of course integrability means that this is a well defined expression,
  we can use any sequence of partitions of $[0,1]$ and any choice of
  points in the partitions and get the same limit.
  In particular,
  if we choose $\set{[0,\frac{1}{k},\dots,\frac{k-1}{k},1]}_{k=1}^{\infty}$
  with points $\set{\xi^k_j := \frac{j}{k}}_{j=1}^{k}$ for the $k$th partition,
  we get that:
  \[
    \E[H(X)] =
    \int_{0}^{1} H(x) \dx =
    \lim_{n \to \infty} \frac{1}{n} \sum_{j=1}^{n} H\del{\frac{j}{n}} =
    \lim_{n \to \infty} \E\left[H(X_n)\right],
  \]
  which completes the proof.
\end{solution}

\newpage

\section{Random Vectors}

\begin{exercise}
  Let $(X,Y)$ be a point chosen uniformly at random from the unit square
  $[0,1] \times [0,1]$.
  \begin{enumerate}
    \item[(1)] Show that $X$ and $Y$ are independent.
    \item[(2)] Let $D$ denote the distance between $X$ and $Y$ in part $(1)$.
      Calculate $\E[D^2]$.
      (Hint: You don't have to calculate the density function of $D$ or $D^2$).
  \end{enumerate}
\end{exercise}
\begin{solution} \phantom{}
  \begin{enumerate}
    \item[(1)] The joint density of $(X,Y)$ is:
      \[
        f_{(X,Y)}(x,y) :=
        \begin{cases}
          1 &\text{if $0 \le x \le 1$ and $0 \le y \le 1$}, \\
          0 &\text{otherwise}.
        \end{cases}
      \]
      We can compute the marginal densities:
      \begin{align*}
        f_X(x) &=
        \int_{0}^{1} f_{(X,Y)}(x,y) \dy =
        \int_{0}^{1} 1 \dy = 1 \\
        f_y(y) &=
        \int_{0}^{1} f_{(X,Y)}(x,y) \dx =
        \int_{0}^{1} 1 \dx = 1.
      \end{align*}
      Since we have and equality:
      \[
        f_{(X,Y)}(x,y) = 1 = f_X(x) f_y(y),
      \]
      it follows by a theorem from class that $X$ and $Y$ are independent.
    \item[(2)] We have that $D = |X - Y|$ and thus $D^2 = (X - Y)^2$.
      By linearity of expectation:
      \[
        \E[D^2] =
        \E[(X-Y)^2] =
        \E[X^2] - 2\E[XY] + \E[Y^2].
      \]
      Since by $(1)$ we know that $X$ and $Y$ are independent
      and $\E[X] = \E[Y] = \int_{0}^{1} x \dx = \half$ we can compute that:
      \[
        \E[XY] = \E[X] \E[Y] = \half \cdot \half = \frac{1}{4}.
      \]
      We can also use the law of the unconscious statistician in order
      to get:
      \[
        \E[X^2] = \int_{0}^{1} x^2 \dx = \frac{1}{3}.
      \]
      By symmetry, the same holds for $Y$ and so $\E[Y^2] = \frac{1}{3}$.
      In conclusion,
      \[
        \E[D^2] =
        \E[X^2] - 2\E[XY] + \E[Y^2] =
        \frac{1}{3} - 2 \cdot \frac{1}{4} + \frac{1}{3} =
        \frac{1}{6},
      \]
      which completes this exercise.
   \end{enumerate}
\end{solution}

\begin{exercise}
  Suppose that $n$ balls are placed uniformly at random in $r$ boxes
  (a box can contain more than one ball).
  Let $X_i$ equal $1$ if box $i$ is empty and equal $0$ otherwise.
  Let $N$ denote the number of empty boxes.
  \begin{enumerate}
    \item[(1)] Calculate $\E[X_i]$.
    \item[(2)] Calculate $\E[X_i X_j]$.
    \item[(3)] Calculate $\E[N]$ and $\sigma^2(N)$.
      (Hint: Represent $N$ as $N = \sum_{i=1}^{r} X_i$).
  \end{enumerate}
\end{exercise}
\begin{solution} \phantom{}
  \begin{enumerate}
    \item[(1)] Note that $X_i$ are discrete and thus:
      \[
        \E[X_i] :=
        \sum_{a \in \set{0,1}} a \cdot \Prob(X_i = a) =
        \Prob(X_i = 1) =
        \del{\frac{r - 1}{r}}^n =
        \del{1 - \frac{1}{r}}^n,
      \]
      since the probability that the $i$th box is empty is exactly
      the probability of all balls going to the other boxes,
      and since they are placed at the boxes uniformly and independently
      we get this result.
    \item[(2)] Note that $X_i X_j = X_i$ for $i = j$
      and thus we can use the result from $(1)$.
      However when $i \neq j$, we get that $X_i X_j$ is an indicator
      function that is $1$ if and only if both the $i$th and the $j$th
      boxes are empty.
      This implies that:
      \[
        \E[X_i X_j] = \Prob(X_i X_j = 1) = \Prob(\text{$i$th and $j$th boxes
        empty}) = 
        \del{\frac{r - 2}{r}}^n =
        \del{1 - \frac{2}{r}}^n,
      \]
      since the probability of each ball to land outside these two boxes
      is $\frac{r - 2}{r}$ (they are placed uniformly) and the placing
      of the balls is independent.
    \item[(3)] We have that $N = \sum_{i=1}^{r} X_i$
      and we need to find $\E[N]$.
      From linearity of expectation and $(2)$ we then see that:
      \[
        \E[N] =
        \E\left[\sum_{i=1}^{r} X_i\right] =
        \sum_{i=1}^{r} \E[X_i] = 
        r \cdot \del{\frac{r - 1}{r}}^n =
        \frac{(r - 1)^n}{r^{n-1}}.
      \]
      We are now tasked with finding the variance of $N$
      denoted $\sigma^2(N)$ or $\Var(N)$.
      We know that:
      \[
        \Var(N) =
        \E[N^2] - \E[N]^2,
      \]
      and by linearity of expectation we have:
      \[
        \E[N^2] =
        \E\left[\del{\sum_{i=1}^{r} X_i}^2\right] =
        \sum_{i=1}^{r} \E[X_i^2] + 2 \cdot \sum_{i < j} \E[X_i X_j],
      \]
      and also:
      \[
        \E[N]^2 =
        \E\left[\sum_{i=1}^{r} X_i\right]^2 =
        (\sum_{i=1}^{r} \E[X_i])^2 =
        \sum_{i=1}^{r} \E[X_i]^2 + 2 \cdot \sum_{i < j} \E[X_i] \E[X_j],
      \]
      so we conclude that:
      \[
        \Var(N) =
        \sum_{i=1}^{r} (\E[X_i^2] - \E[X_i]^2) +
        2 \cdot \sum_{i < j} (\E[X_i X_j] - \E[X_i] \E[X_j]),
      \]
      and by plugging in $\Var(X_i) = \E(E_i^2) - \E[X_i]^2$
      and $\Cov(X_i,X_j) = \E[X_i X_j] - \E[X_i] \E[X_j]$ we get:
      \[
        \Var(N) = 
        \sum_{i=1}^{r} \Var(X_i) +
        2 \cdot \sum_{i < j} \Cov(X_i,X_j),
      \]
      and we can compute these from the same formulas from earlier.
      Notice that we have $X_i^2 = X_i$ and so from $(1)$ we get:
      \[
        \Var(X_i) =
        \E[X_i^2] - \E[X_i]^2 =
        \E[X_i] - \E[X_i]^2 =
        \del{\frac{r - 1}{r}}^n \del{1 - \del{\frac{r - 1}{r}}^n}.
      \]
      From parts $(1)$ and $(2)$ we get:
      \[
        \Cov(X_i,X_j) =
        \E[X_i X_j] - \E[X_i] \E[X_j] =
        \del{\frac{r - 2}{r}}^n - \del{\frac{r - 1}{r}}^{2n}.
      \]
      In conclusion since we have $\binom{r}{2}$ pairs $(i,j)$ with $i < j$
      in the sum, we have that:
      \begin{align*}
        \Var(N) &= 
        \frac{(r - 1)^n}{r^{n-1}} \del{1 - \del{\frac{r - 1}{r}}^n} +
        2 \cdot \binom{r}{2} \cdot 
        \del{\del{\frac{r - 2}{r}}^n - \del{\frac{r - 1}{r}}^{2n}} \\ &=
        \frac{(r - 1)^n}{r^{n-1}} \del{1 - \del{\frac{r - 1}{r}}^n} +
        r (r - 1) \cdot 
        \del{\del{\frac{r - 2}{r}}^n - \del{\frac{r - 1}{r}}^{2n}}.
      \end{align*}
  \end{enumerate}
\end{solution}

\begin{exercise}
  Let $X$ and $Y$ be i.i.d. random variables distributed
  according to $\Exp(\lambda)$.
  Let $Z = \frac{Y}{X}$.
  Find the distribution function $F_Z(z)$ and the density function $f_Z(x)$.
\end{exercise}
\begin{solution}
  Let $z > 0$.
  We have:
  \[
    F_Z(z) =
    \Prob(Z \le z) =
    \Prob\left(\frac{Y}{X} \le z\right) =
    \Prob\del{(X,Y) \in A_z} =
    \int_{A_z} f_{(X,Y)}(x,y) \dx \dy,
  \]
  where $A_z := \set{(x,y) \mid x > 0 \stand y > 0 \stand y < zx}$,
  and $f_{(X,Y)}(x,y)$ is the joint density function of $(X,Y)$.
  Since $X$ and $Y$ are i.i.d we have $f_{(X,Y)}(x,y) = f_X(x) f_Y(x)$,
  and since they both distribute according to $\Exp(\lambda)$,
  we have that $f_{(X,Y)}(x,y) = \lambda^2 e^{-\lambda(x + y)}$.
  We now compute the integral using Fubini's theorem (which is allowed
  since the integrand is nonnegative):
  \[
    F_Z(z) =
    \int_{A_z} \lambda^2 e^{-\lambda(x + y)} \dx \dy =
    \int_{0}^{\infty} \del{\int_{0}^{zx} \lambda^2 e^{-\lambda(x + y)} \dy} \dx.
  \]
  We can compute the inner integrand easily:
  \[
    \int_{0}^{zx} \lambda^2 e^{-\lambda(x + y)} \dy =
    \lambda^2 e^{-\lambda x} \int_{0}^{zx}  e^{-\lambda y)} \dy =
    \lambda e^{-\lambda x} (1 - e^{-\lambda z x}),
  \]
  where the last step follows from the fundamental theorem of calculus.
  We now have:
  \[
    F_Z(z) =
    \int_{0}^{\infty} \lambda e^{-\lambda x} (1 - e^{-\lambda z x}) \dx =
    \int_{0}^{\infty} \lambda e^{-\lambda x} \dx
      - \int_{0}^{\infty} \lambda e^{-\lambda x(z + 1)} \dx.
  \]
  The first integral is:
  \[
    \int_{0}^{\infty} \lambda e^{-\lambda x} \dx =
    [-e^{-\lambda x}]\biggr\vert^{\infty}_{0} = [0 - (-1)] = 1,
  \]
  and the second integral is exactly the first integral after a change of
  variables $x \mapsto (z+1)x$, and since the derivative of this map
  is $(z + 1)$ we get that the second integral is equal to $\frac{1}{z + 1}$.
  This implies that:
  \[
    F_Z(z) = 1 - \frac{1}{z+1} = \frac{z}{z + 1},
  \]
  for $z > 0$ (recall that in the beginning we assumed $z > 0$).
  Since $X \sim Y \sim \Exp(\lambda)$, both of these variables are positive,
  and thus for $z \le 0$ we have $F_Z(z) = 0$.
  We conclude that:
  \[
    F_Z(z) =
    \begin{cases}
      \frac{z}{z + 1} &z > 0, \\
      0 &z \le 0.
    \end{cases}
  \]
  In order to find the density function of $Z$,
  we simply differentiate $F_Z(z)$ and get that:
  \[
    f_Z(z) = \frac{\dif F_Z(z)}{\dz} =
    \begin{cases}
      \frac{1}{(z + 1)^2} &z > 0, \\
      0 &z \le 0.
    \end{cases}
  \]
  This computation completes the proof.
\end{solution}

\begin{exercise}
  Consider tossing a coin $n$ times.
  We define a run to be a sequence of tosses that result in the same outcome.
  Thus, for example, if $n = 7$, the sequence $HHTHTTH$ contains 5 runs,
  while if $n = 12$ the sequence $HHHHHTTTTHTT$ contains $4$ runs.
  Fix $n > 0$ and let $p$ be the probability of heads.
  Let $R$ denote the number of runs.
  Calculate $\E[R]$, the expected value of the number of runs.
  (Hint: Use the method of indicator random variables).
\end{exercise}
\begin{solution}
  Define the indicator $I_k$ by:
  \[
    I_k :=
    \begin{cases}
      1 &\text{a new run starts at the $k$th toss}, \\
      0 &\text{otherwise},
    \end{cases}
  \]
  for $1 \le n$.
  Note that $I_1 = 1$ always and for $k \geq 2$ we have $I_k = 1$ if 
  and only if the $k$th toss is different than the $k-1$th toss.
  Also note that we have $R = \sum_{k=1}^{n} I_n$.
  We can now compute the expectation of the indicators
  using the previous conclusions and get $\E[I_1] = 1$
  and for $k \geq 2$ get:
  \begin{align*}
    \E[I_k] &=
    \Prob(\text{$k$th toss is different than the $k-1$th toss}) =
    \Prob(HT) + \Prob(TH) \\ &= (1 - p)p + p(1 - p) = 2p(1 - p),
  \end{align*}
  where the second equality follows from the fact we only have two options
  for the event to happen; either $H$ then $T$ or $T$ then $H$ must happen.
  We compute these probabilities easily since the tosses are independent.
  Thus, by linearity of expectation we get:
  \[
    \E[R] =
    \E\left[\sum_{k=1}^{n} I_n\right] =
    \sum_{k=1}^{n} \E[I_n] =
    1 + (n - 1) 2p (1 - p),
  \]
  which completes the proof.
\end{solution}

\newpage

\section{Conditional Distributions}

\begin{exercise}
  Let $n \in \N$.
  Let $(X,Y)$ be a discrete random vector
  with $p(x,y) = \binom{n}{y} p^{y+1} (1 - p)^{n - y + x}$,
  for $y = 0,1,\dots,n$ and $x = 0,1,2,\dots$.
  \begin{enumerate}
    \item[(1)] Find $p_X(x)$ and $p_{Y \mid X}(y \mid x)$.
    \item[(2)] Are $X$ and $Y$ independent?
    \item[(3)] Find $\E[XY]$ without doing a calculation,
      just use results you already know.
  \end{enumerate}
\end{exercise}
\begin{solution} \phantom{}
  \begin{enumerate}
    \item[(1)] In order to find the marginal distribution of $x$ we simply 
      calculate by the formula:
      \[
        p_X(x) =
        \sum_{y=0}^{n} \binom{n}{y} p^{y+1} (1-p)^{n-y+x} =
        p (1-p)^x \sum_{y=0}^{n} \binom{n}{y} p^{y} (1-p)^{n-y} =
        p (1-p)^x,
      \]
      since the sum is the binomial expansion of $(p + (1 - p))^n$ which
      is simply $1$.  
      As a corollary we get that $X \sim \Geo(p)$ on $0,1,2,\dots$.

      In order to find $p_{Y \mid X}(y \mid x)$ we are still going
      to calculate by the formula:
      \[
        p_{Y \mid X}(y \mid x) =
        \frac{p(x,y)}{p_X(x)} =
        \frac{\binom{n}{y} p^{y+1} (1 - p)^{n - y + x}}{p (1-p)^x} =
        \binom{n}{y} p^y (1 - p)^{n - y}.
      \]
      Therefore we can conclude that $Y \mid X = x \sim \Bin(n,p)$,
      which is independent of $x$ (we will use this in the following part).
    \item[(2)] Indeed, both variables are independent,
      and we can see that by noticing that their distributions factorizes
      as follows:
      \[
        p(x,y) =
        \underbrace{p (1 - p)^x}_{p_X(x)}
        \underbrace{\binom{n}{y} p^y (1 - p)^{n - y}}_{p_Y(y)}.
      \]
    \item[(3)] Since $X$ and $Y$ are independent we get
      $\E[XY] = \E[X]\E[Y]$, and since we have seen that
      $Y \mid X = x \sim \Bin(n,p)$ for all $x$ so 
      $Y \sim \Bin(n,p)$ and $X \sim \Geo(p)$ with support $0,1,2,\dots$
      which have known expectations, we conclude that:
      \[
        \E[XY] =
        \E[X] \E[Y] =
        \del{\frac{1-p}{p}} \del{np} =
        n (1 - p).
      \]
  \end{enumerate}
\end{solution}

\begin{exercise}
  Let $(X,Y)$ be a random vector with
  joint density $f(x,y) = \alpha^2 e^{-\alpha y}$,
  for $0 \le x \le y < \infty$, where $\alpha > 0$.
  \begin{enumerate}
    \item[(1)] Find $f_X(x)$.
    \item[(2)] Find $f_{Y \mid X}(y \mid x)$.
    \item[(3)] Using parts $(1)$ and $(2)$,
      and considering the results we obtained in class for the waiting time
      problem, write down (with a short explanation but no calculation)
      the density of $f_Y(y)$ of $Y$.
  \end{enumerate}
\end{exercise}
\begin{solution} \phantom{}
  \begin{enumerate}
    \item[(1)] We are given that the support is $0 \le x \le y < \infty$,
      so outside of this region the density is zero. 
      However, for $x \geq 0$ we get:
      \[
        f_X(x) =
        \int_{x}^{\infty} \alpha^2 e^{-\alpha y} \dy =
        \alpha^2 \frac{e^{-\alpha x}}{\alpha} =
        \alpha e^{-\alpha x}.
      \]
      We can conclude that $X \sim \Exp(\alpha)$.
    \item[(2)] By definition we get that:
      \[
        f_{Y \mid X}(y \mid x) =
        \frac{p(x,y)}{p_X(x)} =
        \frac{\alpha^2 e^{-\alpha y}}{\alpha e^{-\alpha x}} =
        \alpha e^{-\alpha (y - x)}.
      \]
      Thus, we can conclude that $Y - X \mid X = x \sim \Exp(\alpha)$,
      which is independent of $X$.
    \item[(3)] It is obvious that $Y = X + (Y - X)$,
      and from $(1)$ and $(2)$ we know that $X \sim \Exp(\alpha)$
      and $Y - X \sim \Exp(\alpha)$.
      Thus, by a result from class we get that:
      \[
        f_Y(y) =
        \alpha^2 y e^{- \alpha y}.
      \]
  \end{enumerate}
\end{solution}

\begin{exercise}
  Let $X$ and $Y$ be i.i.d. and distributed uniformly on $[0,1]$.
  Let $Z = X + Y$.
  Find the density $f_Z(x)$ of $Z$.
  (Warning: be careful with your use of convolution!).
\end{exercise}
\begin{solution}
  We have that $f_X(x) = f_Y(x) = \ind_{[0,1]}(x)$,
  and so by definition we get that:
  \[
    f_Z(z) =
    \int_{\R} f_X(t) f_Y(z - t) \dt =
    \int_{\R} \ind_{[0,1]}(t) \ind_{[0,1]}(z - t) \dt,
  \]
  so $f_Z(z)$ is simply the length of the intersection of
  the intervals $[0,1] \cap [z - 1,z]$.
  In case $0 \le z \le 1$ it is clear that the length is $z$
  so $f_Z(z) = 1$.
  In case $1 < z \le 2$ we have that
  $[0,1] \cap [z - 1, z] = [z - 1, 1]$, which has length $2 - z$.
  In other cases the intersection of the intervals is empty and thus
  $f_Z(z) = 0$ so we conclude that:
  \[
    f_Z(z) :=
    \begin{cases}
      z &0 \le 1 \le 1, \\
      2 - z &1 < z \le 2, \\
      0, &\text{otherwise}.
    \end{cases}
  \]
\end{solution}

\begin{exercise}
  Let $(X,Y)$ be a random vector with joint density $f(x,y) = x e^{-x(y + 1)}$,
  for $x,y \geq 0$.
  \begin{enumerate}
    \item[(1)] Find $f_X(x)$ the density of $X$.
    \item[(2)] Find $f_{Y \mid X}(y \mid x)$.
    \item[(3)] Calculate $\E[Y \mid X = x]$.
    \item[(4)] Use the law of total expectation along with $(1)$
      and $(3)$ to show that $\E[Y] = \infty$.
    \item[(5)] Calculate $f_Y(y)$, the density of $Y$ and use it to calculate
      $\E[Y]$ directly.
  \end{enumerate}
\end{exercise}
\begin{solution} \phantom{}
  \begin{enumerate}
    \item[(1)] We compute the density of $X$ by definition:
      \[
        f_X(x) =
        \int_{0}^{\infty} x e^{-x (y + 1)} \dy =
        x e^{-x} \int_{0}^{\infty} e^{-xy} \dy =
        x e^{-x} \cdot \frac{1}{x} =
        e^{-x},
      \]
      for $x \geq 0$, and therefore $X \sim \Exp(1)$.
    \item[(2)] By definition we get:
      \[
        f_{Y \mid X}(y \mid x) =
        \frac{f(x,y)}{f_X(x)} =
        \frac{x e^{-x(y + 1)}}{e^{-x}} =
        x e^{-xy},
      \]
      for $y \geq 0$. 
      Thus we conclude $Y \mid X = x \sim \Exp(x)$.
    \item[(3)] We want to find $\E[Y \mid X = x]$
      which is the expectation of a random variable which we found
      distributes according to $\Exp(x)$ in part $(2)$
      so it is simply $\E[Y \mid X = x] = \frac{1}{x}$.
    \item[(4)] By the law of total expectation,
      and the results from part $(3)$ we get:
      \[
        \E[Y] =
        \E\left[\E[Y \mid X]\right] =
        \E\left[\frac{1}{X}\right].
      \]
      Using the result from part $(1)$ that $f_X(x) = e^{-x}$
      and the LOTUS formula we get:
      \[
        \E\left[\frac{1}{X}\right] =
        \int_{0}^{\infty} \frac{1}{x} e^{-x} \dx = \infty,
      \]
      where the last equality follows from the fact the integral diverges
      as $0$.
      One can see that the integral diverges since the integrand is positive,
      and if we assume by contradiction that it converges by the comparison
      test the integral $\int_{0}^{\infty} \frac{1}{x} e^{-1} \dx$ would
      also diverge, and so will the integral 
      $\int_{0}^{\infty} \frac{1}{x} \dx$ which is a contradiction.
      Hence we conclude:
      \[
        \E[Y] =
        \E\left[\frac{1}{X}\right] = \infty.
      \]
    \item[(5)] We compute $f_Y(y)$ by definition:
      \[
        f_Y(y) =
        \int_{0}^{\infty} x e^{-x(y+1)} \dx =
        \int_{0}^{\infty} x e^{-(y+1)x} \dx =
        \frac{1}{(y + 1)^2},
      \]
      for $y \geq 0$.
      The integral is a famous integral from calculus $2$
      which is solved by the substitution $t = -(y+1)$,
      so we are just using a known result.
      In order to find the expectation by definition we use
      the substitution $u = y + 1$ and compute:
      \[
        \E[Y] =
        \int_{0}^{\infty} \frac{y}{(y + 1)^2} \dy =
        \int_{1}^{\infty} \frac{u - 1}{u^2} \dy =
        \int_{1}^{\infty} \frac{1}{u} - \frac{1}{u^2} \dy = \infty.
      \]
  \end{enumerate}
\end{solution}

\begin{exercise}[Wald's equation]
  Let $\{X_n\}_{n=1}^{\infty}$ be a sequence of i.i.d. random variables
  with finite expectation.
  Denote $\mu = \E[X_n]$.
  Let $N$ be a random variable with finite expectation,
  independent of $\{X_n\}_{n=1}^{\infty}$,
  and distributed on $\N$.
  Let $Z = S_N$,
  that is, $Z$ is a random variable satisfying $Z \mid \set{N = n} \sim S_n$.
  Use the law of total expectation to calculate $\E[Z] = \E[S_N]$.
\end{exercise}
\begin{solution}
  By definition $Z = S_N = \sum_{k=1}^{N} X_k$. 
  We have that $Z \mid \set{N = n} \sim S_n$ and thus:
  \[
    \E[Z \mid N = n] =
    \E[S_n \mid N = n] =
    \E\left[\sum_{k=1}^{n} X_k \mid N = n\right].
  \]
  By linearity of expectation, and since each $X_k \mid N = n$
  has expectation $\E[X_k] = \mu$ (this follows from the fact $N$ is independent
  of $\{X_n\}_{n=1}^{\infty}$)
  we conclude:
  \[
    \E[S_n] =
    \sum_{k=1}^{n} \E[X_k \mid N = n] = n\mu \implies
    \E[Z \mid N = n] = n \mu.
  \]
  We can now apply the law of total expectation on $Z$ and get:
  \[
    \E[Z] =
    \E\left[\E[Z \mid N]\right] =
    \E[N\mu] =
    \mu \E[N].
  \]
\end{solution}

\newpage

\section{Normal Distributions and Random walks}

\begin{exercise}
  A polling organization samples $1200$ voters in a large city in order to
  estimate the proportion of people there who support candidate $A$ in a certain
  election. 
  At least how large (approximately) would the true proportion $p$ of
  people in the city who support $A$ have to be in order that there be at least
  a $95\%$ chance that at least half the people sampled will support $A$?
\end{exercise}
\begin{solution}
  If we denote the number of citizens who voted for $A$ by $X$
  then it is clear $X \sim \Bin(1200,p)$,
  and the exercise reduces to finding approximately minimal $p$
  so that $\Prob(X \geq 600) \geq 0.95$.
  We note that $X = \sum_{i=1}^{1200} X_i$ for random variables $X_i$
  which represent if the $i$th citizen voted for $A$ or not,
  so this problem is a standard problem about the central limit theorem.
  By the central limit theorem we get that $X \approx N(\mu,\sigma^2)$
  where $\mu = 1200p$ and $\sigma = \sqrt{1200 p (1 - p)}$.
  We would like to find $\Pr(X \geq 600)$,
  but recall that the normal distribution is continuous
  and so if we denote $Y \sim N(\mu,\sigma^2) \approx X$
  for the values $\mu$ and $\sigma$ defined earlier,
  a better approximation is given by $\Pr(Y \geq 599.5) \geq 0.95$.
  If we denote $Z = \frac{X - \mu}{\sigma} \sim N(0,1)$,
  then by a property of the normal distribution we get:
  \[
    \Pr(Y \geq 599.5) =
    \Pr\del{Z \geq \frac{599.5 - \mu}{\sigma}}.
  \]
  From standard normal distribution tables
  we get that $\Prob(Z \geq -1.645) \approx 0.95$,
  and therefore we have that:
  \[
    \Pr\del{Z \geq \frac{599.5 - \mu}{\sigma}} \geq 0.95 \iff
    \frac{599.5 - \mu}{\sigma} \le -1.645.
  \]
  By substituting the values of $\mu$ and $\sigma$ we get the inequality:
  \[
    1200p - 599.5 \geq 1.645 \sqrt{1200 p (1 - p)},
  \]
  and by solving this we get that the minimal $p$
  so that $\Prob(X \geq 600) \geq 0.95$ is approximately $0.523$.
\end{solution}

\begin{exercise}
  \label{ex:six-two}
  In this exercise, you will prove that if $X \sim N(\mu_1,\sigma_1^2)$
  and $Y \sim N(\mu_2,\sigma_2^2)$, and $X$ and $Y$ are independent,
  then $X + Y \sim N(\mu_1 + \mu_2, \sigma_1^2 + \sigma_2^2)$.
  \begin{enumerate}
    \item[(1)] Prove the exercise in the case that $\mu_1 = \mu_2 = 0$. 
      Hint: use convolution:
      $f_{X+Y}(z) = \int_{-\infty}^{\infty} f_X(x) f_Y(z - x) \dx$.
      The expression in the exponent appearing in $f_X(x) f_Y(z - x)$
      is $-(\frac{x^2}{2\sigma_1^2} + \frac{(z - x)^2}{2\sigma_2^2})$. 
      Complete the square to rewrite this in the form
      $-(\frac{\sigma_1^2 + \sigma_2^2}{2\sigma_1^2 \sigma_2^2}(x - A)^2 + B)$,
      where $A$ and $B$ depend on $\sigma_1$, $\sigma_2$ and $z$. 
      Now integrate.
    \item[(2)] If $W$ is a random variable with desndity $f(w)$ and $c$
      is a constant, then show
      that the random variable $W + c$ has density $f(w - c)$.
    \item[(3)] Use parts $(1)$ and $(2)$ in order to prove the claim for general
      $\mu_1$ and $\mu_2$. 
  \end{enumerate}
\end{exercise}
\begin{solution} \phantom{}
  \begin{enumerate}
    \item[(1)] As mentioned in the hint,
      we know the densities of these functions,
      and we also know that the density function for the sum $X + Y$
      is given by convolution of each of their density functions
      (the formula works because they are independent).
      Using the expression given in the hint we get that:
      \[
        f_{X+Y}(z) =
        \int_{-\infty}^{\infty} f_{X}(x) f_{Y}(z-x) \dx =
        \frac{1}{2 \pi \sigma_1 \sigma_2} \int_{-\infty}^{\infty}
        \exp\del{-\frac{x^2}{2\sigma_1^2} - \frac{(z - x)^2}{2\sigma_2^2}} \dx.
      \]
      Where $\exp(y)$ is the exponential function $e^y$.
      Ignoring the constant in the exponent we get that:
      \[
        \frac{x^2}{\sigma_1^2} + \frac{(z - x)^2}{\sigma_2^2} =
        \del{\frac{1}{\sigma_1^2} + \frac{1}{\sigma_2^2}} x^2
          -\frac{2z}{\sigma_2^2} x + \frac{z^2}{\sigma_2^2},
      \]
      so by completing the square we get that:
      \[
        \frac{x^2}{\sigma_1^2} + \frac{(z - x)^2}{\sigma_2^2} =
        \frac{\sigma_1^2 + \sigma_2^2}{\sigma_1^2 \sigma_2^2}
          \del{x - \frac{z \sigma_1^2}{\sigma_1^2 + \sigma_2^2}}^2 +
          \frac{z^2}{\sigma_1^2 + \sigma_2^2}.
      \]
      Now by using this result and taking the exponential constant
      independent of $x$
      outside of the integral we get that:
      \[
        f_{X+Y}(z) =
        \frac{1}{2 \pi \sigma_1 \sigma_2}
          e^{\frac{-z^2}{2(\sigma_1^2 + \sigma_2^2)}}
          \int_{-\infty}^{\infty}
          \exp\del{-\frac{\sigma_1^2 + \sigma_2^2}{2 \sigma_1^2 \sigma_2^2}(x - A)^2} \dx,
      \]
      where $A = \frac{z \sigma_1^2}{\sigma_1^2 + \sigma_2^2}$.
      We can now use the formula of the general Gaussian integral and
      get that:
      \[
        \int_{-\infty}^{\infty} e^{-c(x - A)^2} \dx =
        \sqrt{\frac{\pi}{c}},
      \]
      and in our
      case $c = \frac{\sigma_1^2 + \sigma_2^2}{2 \sigma_1^2 \sigma_2^2}$.
      This means that:
      \[
        f_{X+Y}(z) =
        \frac{1}{2 \pi \sigma_1 \sigma_2}
          e^{\frac{-z^2}{2(\sigma_1^2 + \sigma_2^2)}}
          \sqrt{\frac{\pi}{\frac{\sigma_1^2 + \sigma_2^2}{2 \sigma_1^2 \sigma_2^2}}} =
          \frac{1}{\sqrt{2\pi (\sigma_1^2 + \sigma_2^2)}}
          e^{\frac{-z^2}{2(\sigma_1^2 + \sigma_2^2)}},
      \]
      which implies that $X + Y \sim N(0,\sigma_1^2 + \sigma_2^2)$,
      and completes the first part of the exercise.
    \item[(2)] Note that if $W$ has density $f(w)$ and $c \in \R$
      then:
      \[
        \Prob(W + c \le w) =
        \Prob(W \le w - c),
      \]
      so the density of $W + c$ is clearly $f(w - c)$.
    \item[(3)] In the general case we suppose that $X \sim N(\mu_1,\sigma_1^2)$
      and $Y \sim N(\mu_2,\sigma_2^2)$.
      This implies that $X = \mu_1 + X_0$ and $Y = \mu_2 + Y_0$
      where $X_0 \sim N(0,\sigma_1^2)$ and $Y_0 \sim N(0,\sigma_2^2)$,
      which are independent, so from part $(1)$ we get that
      $X_0 + Y_0 \sim N(0,\sigma_1^2 + \sigma_2^2)$.
      Now we have that:
      \[
        X + Y = (X_0 + Y_0) + (\mu_1 + \mu_2),
      \]
      and by part $(2)$ shifting by $\mu_1 + \mu_2$ gives
      that $X + Y \sim N(\mu_1 + \mu_2, \sigma_1^2 + \sigma_2^2)$,
      which completes the proof.
  \end{enumerate}
\end{solution}

\begin{exercise}
  Let $X$ and $Y$ be independent with $X \sim N(0,9)$ and $Y \sim N(1,16)$. 
  Calculate approximately $\Prob(X \le Y)$.
\end{exercise}
\begin{solution}
  Since $-X \sim N(0,9)$ we can simply use the previous exercise
  (\Cref{ex:six-two}) in order to conclude that:
  \[
    Y + (-X) \sim N(1 - 0, 16 + 9) = N(1,25).
  \]
  Now we can simply apply the basic property of the standard distribution:
  \[
    \Prob(X \le Y) =
    \Prob(Y - X \geq 0) =
    \Prob\del{Z \geq \frac{0 - 1}{5}} =
    \Prob(Z \geq -0.2),
  \]
  for $Z \sim N(0,1)$.
  By using symmetry we can conclude that:
  \[
    \Prob(X \le Y) =
    \Prob(Z \geq -0.2) =
    \Prob(Z \le 0.2) =
    \Phi(0.2) \approx 0.579,
  \]
  which concludes the proof.
\end{solution}

\begin{exercise}
  Let $p \in (\half,1)$ and consider the simple random walk satisfying
  $\Prob(X_{n+1} = X_n + 1) = p$ and $\Prob(X_{n+1} = X_n - 1) = 1 - p$. 
  Let $T_y := \inf\set{n \geq 0 \mid X_n = y}$. 
  In class we proved that:
  \[
    \Prob_k(T_0 < \infty) =
    \del{\frac{1 - p}{p}}^k,
  \]
  for $k \geq 0$.
  Now let $M := \inf\set{X_n \mid n \geq 0}$.
  The random variable $M$ is the minimum value of the random walk path.
  \begin{enumerate}
    \item[(1)] Calculate $\Prob_0(M \le y)$, for $y \le 0$.
      Hint: $\set{M \le y} = \set{T_y < \infty}$.
    \item[(2)] Calculate $\Prob_0(M = y)$, for $y \le 0$.
    \item[(3)] Calculate $\E_0[M]$.
  \end{enumerate}
\end{exercise}
\begin{solution} \phantom{}
  \begin{enumerate}
    \item[(1)]
      For $y \le 0$, we notice that $\set{M \le y} = \set{T_y < \infty}$.
      We also note that by symmetry starting from $0$ and hitting $y \le 0$
      is the same is starting from $-y \geq 0$ and hitting $0$.
      By the given result from class,
      \[
        \Prob_0(T_y < \infty) =
        \Prob_{-y}(T_0 < \infty) = \del{\frac{1 - p}{p}}^{-y},
      \]
      which completes the first part of the exercise.
    \item[(2)] For convenience, from now on we denote
      \[
        r := \frac{1 - p}{p}.
      \]
      In this part of the exercise we just need to use differences
      since the probability of $M = y$ is exactly the probability
      of $\set{M \le y}$ minus the probability of $\set{M \le y-1}$.
      From $(1)$ we get that:
      \[
        \Prob_0(M = y) =
        \Prob_0(M \le y) - \Prob_0(M \le y - 1) =
        r^{-y} - r^{-(y - 1)} =
        (1 - r) r^{-y},
      \]
      which completes the proof of the second exercise.
    \item[(3)] Notice that $\Prob_0(M \le 0) = 1$
      since the first step is $0$,
      so $\Prob_0(M = y) = 0$ for each $y > 0$.
      Thus, we have by definition and part $(2)$ that:
      \[
        \E_0[M] =
        \sum_{y \le 0} y \Prob_0(M = y) =
        (1 - r) \sum_{y \le 0} y r^{-y} =
        (1 - r) \sum_{k=1}^{\infty} (-k) r^k.
      \]
      Recall that $\sum_{k=0}^{\infty} r^k = \frac{1}{1 - r}$,
      and therefore by termwise differentiation we get
      that
      \[
        \sum_{k=1}^{\infty} k r^{k-1} = \frac{1}{(1 - r)^2}.
      \]
      Then by multiplying the series by $(-1)$ and $r$
      we get that $\sum_{k=1}^{\infty} (-k) r^{k} = -\frac{r}{(1 - r)^2}$,
      and thus we conclude that:
      \[
        \E_0[M] =
        (1 - r) \sum_{k=1}^{\infty} (-k) r^k =
        (1 - r) \del{-\frac{r}{(1 - r)^2}} =
        \frac{r}{r - 1}.
      \]
      Finally, since $r := \frac{1 - p}{p}$, we can write more cleanly:
      \[
        \E_0[M] =
        \frac{r}{r - 1} =
        \frac{1 - p}{1 - 2p}.
      \]
  \end{enumerate}
\end{solution}

\end{document}
